% !TEX root = ../bb_AiRa_Kappaleet.tex

%% HARRIS: Chapter 10

% ö = \%c3\%b6
% ä = \%C3\%A4

\xdef\sectiontitle{\IIsectiontitle} %aiheuttama
\TOCrefs

%%%%%%%%%%%%%%%
\begin{frame}[label=2_1_a]%[label=2_0_TOC1]
\frametitle{\texorpdfstring{\culine[\sectiontitleulinecolor]{\sectiontitle}}{\sectiontitle} }
\label{\TOCname}

\vspace{\chapterTOCvksip}
\begin{itemize}[leftmargin=0.5cm,itemsep=3mm]
    \item[2.1]<1-> \hyperlink{2_1_a}{\CDAlert[blue]{\IIaaSubsectiontitle}}
    \item[2.2]<1-> \hyperlink{2_2_a}{\IIaSubsectiontitle}
    \item[2.3]<1-> \hyperlink{2_3_a}{\IIbSubsectiontitle}	
    \item[2.4]<1-> \hyperlink{2_4_a}{\IIcSubsectiontitle}	
    \item[2.5]<1-> \hyperlink{2_5_a}{\IIdSubsectiontitle}
    \item[2.6]<1-> \hyperlink{2_6_a}{\IIeSubsectiontitle}
    \item[2.7]<1-> \hyperlink{2_7_a}{\IIfSubsectiontitle}
    \item[2.8]<1-> \hyperlink{2_8_a}{\IIgSubsectiontitle}
    \item[2.9]<1-> \hyperlink{2_9_a}{\IIhSubsectiontitle}
    \item[2.10]<1-> \hyperlink{2_10_a}{\IIiSubsectiontitle}
\end{itemize}

%\endofslide 
\end{frame}
%%%%%%%%%%%%%%%

\xdef\sectiontitle{\IIaaSubsectiontitle} 

%%%%%%%%%%%%%%%
\begin{frame}%[label=2_1_a]
\frametitle{\texorpdfstring{\culine[\sectiontitleulinecolor]{\sectiontitle}}{\sectiontitle} }
\framesubtitle{Aineen rakenteen suuruusluokista}
\appear{}

\begin{itemize}[itemsep=4mm]
	\item \CDAlert[red]{Etäisyyksien kirjo} atomien\appear{, niiden ytimien}\appear{, ja molekyylien välillä \CDAlert[red]{on valtava!}}
	\item Esimerkiksi molekyylit ovat usein \appear{kokoluokkaa 0,1–1\;nm:} \\
\begin{itemize} \semismall
	\item typpimolekyyli $d=0,37 \mathrm{\;nm}$
	\item vesimolekyyli $d=0,3 \mathrm{\;nm}$
\end{itemize}

\item Kiinteissä aineissa\appear{/molekyyleissä} \appear{\CDAlert[blue]{sidospituudet} ovat $\sim$0,1–0,2\;nm} \appear{= 1–2\;å}

\item \CDAlert[blue]{Bohrin säde on $a_{0}=0,0529$\;nm} \appear{(vetyatomin perustilan elektronin kiertoradan säde)}

\item \CDAlert[fuchsia]{Atomiytimen} koko on ${\sim}1-10\;\mathrm{fm}=1 ... 10 \cdot 10^{-6}\;\mathrm{nm}$\appear{, mikä on huomattavasti vähemmän kuin atomien tyypilliset sidospituudet}
\implies{ \Alert[tuniblue]{Sidoksien muodostumiseen vaikuttaa suurimmaksi osaksi vain uloimmat elektronit, eli} \CDAlert[tuniblue]{valenssielektronit} }
\end{itemize}


\endofslide 
%\endofsection
\end{frame}
%%%%%%%%%%%%%%%

%%%%%%%%%%%%%%%
\begin{frame}
\frametitle{\texorpdfstring{\culine[\sectiontitleulinecolor]{\sectiontitle}}{\sectiontitle} }
\framesubtitle{Atomien välisten sidosten muodostuminen}
\appear{}

\begin{itemize}
	\item Lähekkäin olevilla atomeilla\appear{, niiden elektronit alkavat nähdä \CDAlert[fuchsia]{'moniatomisen potentiaalienergian'}}
	\appear{\xdef\loc{south east}
\begin{tikzpicture}[remember picture,overlay] % anchor=south
    \node (A) [yshift=-0.25*\figYshift,xshift=\figXshift, anchor=\loc] at (current page.\loc) {\includegraphics[page=1,width=10.5cm]{Figures_booklet/SOM_10_Bonding_figures/SOM_Bonding_Figure_1.png}}; %scale=\figscale. % 8cm max height
\end{tikzpicture}\footnote{\HarrisCR}}
	\item Kuva~1 hahmottelee alimmat energiatilat\appear{/aaltofunktiot} \appear{kahdelle vierekkäiselle ja $L$-levyiselle \CDAlert[blue]{äärelliselle neliökaivolle}} \appear{tilanteissa joissa atomien etäisyys $a$ muuttuu} 
	
\item Kun atomit tulevat lähekkäin\appear{, sekä aaltofunktiot} \appear{että energiatasot muuttuvat.} \appear{Etäisyyden pienentäminen kasvattaa energiatilojen välisiä etäisyyksiä}
\end{itemize}

\endofslide 
%\endofsection
\end{frame}
%%%%%%%%%%%%%%%

%%%%%%%%%%%%%%%
\begin{frame}
\frametitle{\texorpdfstring{\culine[\sectiontitleulinecolor]{\sectiontitle}}{\sectiontitle} }
\framesubtitle{Atomien välisten sidosten muodostuminen}
\appear{}

\begin{itemize}
	\item Jos vierekkäisiä potentiaalikaivoja olisi kolme\appear{, etäisyyksien pienentyessä} \appear{tilalla $n=1$ olisi \Alert[red]{olisi kolme molekyylitilaa} } 
\end{itemize}
% 6.5 cm = midway, 10 cm = 3/4 
\begin{minipage}{6.2cm}
\begin{itemize}
	  \item[] hiukan poikkeavilla energiatasoilla\appear{, jotka etäisyyden kasvaessa muodostaisivat kolme erillistä \\ $n=1$ tilaa}
\item Kun tarkastelu yleistetään \Alert[blue]{$N$:n neliökaivon systeemiin}\appear{, atomien etäisyyksien pienentyessä}\appear{ energiatasot levenevät ja muodostavat \CDAlert[blue]{$N$:n tilan energiavöitä/-}\\ \CDAlert[blue]{kaistoja}}
\end{itemize}
\end{minipage}

\appear{\xdef\loc{south east}%
\begin{tikzpicture}[remember picture,overlay]% anchor=south
    \node (A) [yshift=-0.25*\figYshift,xshift=\figXshift, anchor=\loc] at (current page.\loc) {\includegraphics[page=1,height=7cm]{Figures_booklet/SOM_10_Bonding_figures/SOM_Bonding_Figure_2.png}};%scale=\figscale. % 8cm max height
\end{tikzpicture}%
\footnote{\HarrisCR}}%
%\endofslide 
\endofsection
\end{frame}
%%%%%%%%%%%%%%%

